\section{Method}
\label{sec:method}

\paragraph{Tasks and scoring.}
Study 1's tasks were admitted after a single-agent pilot screened about 45 designs for a discriminating pass rate. The admitted tasks each contain one known trap. \texttt{stamp-interpreter} and \texttt{stamp-2} ask for programs implementing the same invented-language scoping rule; they are two fixed instances of one family, not independent evidence from two domains. \texttt{bench-printf-format} asks for exact formatting where a natural implementation using the host runtime's shortest representation disagrees with the specified output. The same hidden, machine-checked fixtures scored every arm. The seeds vary attempts on fixed task instances, rather than sampling new tasks. The task suite therefore supports comparisons on these traps, not an estimate across software engineering work.


The pilot targeted single-agent pass rates between $0.3$ and $0.7$, inclusive. The interpreter admission pilots used ten seeds each; the retained, raw-backed printf v1 pilot used fifteen seeds; v2 used ten, and five earlier v1 seeds lacked retained artifacts. Two versions of the printf task were piloted, and the one used here was chosen after both pilots' results had been seen, a post-data selection.

\paragraph{Study 1 design.}
For each task, A and C ran on seeds 101--140, interleaved per seed. A's budget ceiling was tied to C's realized spend; matching did not force the same realized spend. K was separately pre-registered after A and C results but before K ran, on the same seeds and tasks. Its number of attempts was chosen from the earlier mean A and C costs, with each attempt run afresh. The selector had no oracle input. The fixed K group sizes were ten attempts for \texttt{stamp-interpreter}, eight for \texttt{stamp-2}, and seven for printf. The per-attempt cap left room for an ordinary A-style attempt; total group cost included every attempt even if it was killed or not selected. The matching target was mean realized C spend, not identical spend on every seed.

K ran from 13:22 to 15:33 UTC, before the account switch recorded at 16:38 UTC; an earlier shift-time estimate was superseded by that switch log. The comparison between K and C still spans separate collection windows. K exceeded its paired C spend on 19, 10, and 9 of the 40 seeds for the three tasks, despite matching at the mean.

The original A--C protocol specified a two-proportion z-test, using Fisher's exact test for small cells, separately by task. After the data, the two interpreter tasks were pooled as a family and A--C was retested with Fisher's exact test and Holm correction across the two families. We report that revised family result alongside the original per-task tests. The K pre-registration fixed six two-sided Fisher comparisons, K against A and C on each task, with Holm correction across all six. Pooled K comparisons on the interpreter family are exploratory and do not replace those registered tests.


The common seed labels put arms on the same task brief, but the attempts are separate model runs; the Fisher calculations do not treat them as paired observations.

\paragraph{Study 2 design and missingness.}
Study 2 specified five arms on seeds 501--520: A, AH, B, K and C. Its primary tests used the short-thinking stratum. The account switcher was not known to the harness at launch. A sanitized log later associated each run with its active account; mixed-account seeds were separated from within-regime comparisons. The run stopped at 190/200 cells, leaving ten absent. Printf seed 501 is kept in descriptive totals but excluded from inference after a documented cap deviation. One printf chatroom run has unknown cost. The long-thinking tests were added after much of that stratum had been observed, and use 19 seeds per arm on \texttt{stamp-interpreter} and 15 on printf. They use Fisher tests with Holm correction across the registered comparison family of fourteen comparisons, but their timing makes them exploratory. Strata condition on A's realized post-run thinking, rather than an independent pre-run measurement; its sentinel also cannot rule out a change later in the same seed.

\paragraph{Account measurement.}
The account analysis joins single-seat run intervals to switch times and uses only runs with thinking telemetry for the thinking-token comparison. Long thinking was defined in the Study 2 pre-registration as at least 4,000 thinking tokens. Five of 1,947 eligible runs straddled a switch and were excluded. Output-token distributions give a second descriptive account comparison, but output across a three-seat room is not directly comparable to a single seat. Neither the active account nor effort was randomized. The client records the association, not the server-side reason for it.

\paragraph{Study 3 design.}
Study 3 asked a different question from Studies 1 and 2: given 20 independent, testable items from one repository and 30 minutes, how many items does each setup land? The setups were one seat with all 20 items (Solo); three seats each given a fixed third of the items and no chat, whose branches the harness merges afterwards in a fixed order, keeping the earlier branch's version on a conflict (Split); a room of three workers and a verifier that organises itself from the list and requires verification (Room3); and the same room with fourteen workers (Room15). Every seat ran \texttt{claude-opus-5-5} from the brief alone, in its own git worktree of a fresh repository holding only the pool's base tree, and was told the deadline. Rooms were scored on the integration branch they declared.

Each item was written by an AI curator as a hidden acceptance test plus a reference fix, and entered a pool only if its test failed three times out of three at the base and passed three times out of three with the fix. Briefs state the symptom, the required behaviour and the interface the test checks, never the cause, the files to change or the approach; from the second pool on, a leak reviewer compared each brief with its fix and a blind fairness reviewer, who never saw the fix, checked that the test asked only for what the brief states. No two items' reference fixes change the same function. Three pools were planned in a fixed order: this system (pool 1), a browser game (pool 2) and a larger product repository frozen at one commit, dependencies included (pool 3). After pool 1 reached the ceiling, later items had to take one agent longer: pool 2's replacement items and every pool 3 item needed a reference fix touching at least two source files and adding at least 50 lines, and the project's own test suite had to pass with each fix. Before pools 2 and 3 were locked, a practice Solo run acted as a ceiling gate: a pool locked if it landed 16 or fewer, otherwise its smallest items were swapped for larger reserves and the gate was run once more.

Each setup ran twice per pool, one run at a time, in a seeded random order. The primary measure is the number of items whose hidden test passes at the setup's final head. Secondary measures are cost from the command-line client's own usage records, cost per passing item, items attempted, merge conflicts and whether the project's suite still passes at the final head. Any access to the hidden tests would void a run. The pre-registration fixed three predictions (Split lands more than Solo; Room3 and Split do not differ; Room15 lands more than Room3) and a direction rule: a difference counts when two setups' means differ by 3 or more items, in the same direction in at least two of the three pools. With two repeats per setup, no significance test is run.

